\section{Replacement and functional universality}\label{sec:replacement}

We now transfer the Gaussian limit to Rademacher couplings.  The comparison
uses the spin autocorrelations themselves.  Its principal input is the
averaged clock-score estimate of Proposition~\ref{prop:clock-scores}; the remaining
point is to differentiate every survival and failure probability in the
history construction.  Conditional copies of a checked path provide these
derivatives while preserving the probability laws to which the score
estimates apply.

\begin{proposition}[Transfer to Rademacher disorder]\label{prop:universality}
For the Rademacher coupling law,
\[
 \Law_{\mathrm{Rad}}\bigl(A_{n,J},B_{n,J}\bigr)
 \weakto \Law_g\bigl(A_g,B_g\bigr)
\]
in
$C_{\mathrm{loc}}((0,\infty))\times
C_{\mathrm{loc}}((0,\infty)^2)$, where $(A_g,B_g)$ is the Gaussian limit
obtained by combining Proposition~\ref{prop:gaussian-subsequence} with the
identification of the entrance family in Theorem~\ref{thm:entrance-selection}.
In particular the entrance family and the semigroup in this limit are the same $h_s^g$ and
$S_t^g=\exp(-t c_*H_g)$ as in the Gaussian case, and the finite time scale is
$T=T_n=n^{2/3}$.
\end{proposition}

Throughout this section a length $D$ of a path interval is in site-clock
units.  All query legs have lengths in $[cT,CT]$, where $0<c<C<\infty$ are
fixed by a compact set of query times.  We put $K=24$.  The small parameter
$\beta>0$ and the convention $\lesssim_\beta$ are those of the preceding two
sections.  Constants in an exponent $O(\beta)$ will depend on $K,c,C$, but
not on the number of conditional-copy generations introduced below.

\subsection{Derivatives of tested histories}

We first state the derivative estimate in the form needed for smooth tests
of several traces.  Let $e=\{i,j\}$, $q_e(x)=x_i x_j$,
$\bar q_e=\mu q_e$, and $\bar q=(\bar q_e)_e$.  Add $\vartheta$ to the
interaction coefficient $J_{ij}/\sqrt n$ and set $u=\tanh\vartheta$.
For a checked full history or prefix $\alpha$ of
Definition~\ref{hist:tested-law}, define the unnormalized checked integral
\begin{equation}\label{repl:tested-integral}
 Z_{\alpha,e}(u)=
 \mathbb E_{\alpha,e,u}\left[
   v_\alpha
   \prod_{I\in\mathscr S_\alpha}w_{I,e}(u)
   \prod_{I\in\mathscr F_\alpha}(1-w_{I,e}(u))\right].
\end{equation}
Here $\mathbb E_{\alpha,e,u}$ is ordinary dynamics on the prescribed full
history or prefix, with uniform or perturbed stationary initialization as
specified by $\alpha$.  The set $\mathscr S_\alpha$ consists of its included
survival checks, and $\mathscr F_\alpha$ is empty or consists of its specified
first failure.  The nonnegative function $v_\alpha$ depends only on the
designated initial, leg, and forest endpoints; it is independent of $u$ and
common to every $e$.  Write $z_\alpha=Z_{\alpha,e}(0)$, which does not depend
on $e$.  When $z_\alpha$ satisfies the inverse-polynomial lower bound in
Definition~\ref{hist:tested-law}, normalizing this base integral gives the
tested path law.  A product of such integrals always uses independent histories;
the external tests factor separately across them.

For a vector indexed by $e$, the norms $|\cdot|_p$ use uniform probability
on the $\binom n2$ edges, as in Proposition~\ref{prop:clock-scores}.  Coefficients
$[u^j]$ below are at $u=0$.  The distinction between coefficients and
derivatives costs constants depending only on $K$.

\begin{lemma}[Averaged derivatives of histories]\label{repl:history-derivatives}
There are constants $a>0$ and $C_K<\infty$ with the following property for
all sufficiently small fixed $\beta>0$.  Fix $A,B<\infty$, at most $K$
histories of the form \eqref{repl:tested-integral}, and integers
$j_\alpha\ge0$ with $1\le k=\sum_\alpha j_\alpha\le K$.  Suppose
$0\le v_\alpha\le n^B$.  On the deterministic static support of
Lemma~\ref{hist:static-support}, put
\[
 \mathscr D_e=\prod_\alpha [u^{j_\alpha}]Z_{\alpha,e}(u).
\]
Then, uniformly in the selected histories,
\begin{align}
 |\mathscr D|_1
 &\le C\left(\prod_\alpha z_\alpha\right)
 \begin{cases}
 n^{-a+C_K\beta},&1\le k\le4,\\
 n^{C_K\beta}T^{(k-4)/2},&5\le k\le K,
 \end{cases}
 +n^{-A},                                                     \label{repl:averaged-history}\\
 |\mathscr D|_\infty
 &\le C n^{C_K\beta}T^{k/2}\prod_\alpha z_\alpha+n^{-A}.
                                                               \label{repl:sup-history}
\end{align}
The constant $C$ and the threshold for $n$ may depend on $A,B$ and on the
fixed rectangles.  The exponent constant $C_K$ is independent of these
precision parameters and of $\beta$.  The same estimates hold for specified
derivatives in $\vartheta$ in place of the coefficients in $u$.  Signed tests
are covered by applying the statement to their positive and negative parts.
\end{lemma}

\begin{proof}
The additive precision in the lemma is useful for very small tested
masses.  To prove it, first suppose each $z_\alpha\ge n^{-M}$ for a fixed
$M$.  Normalize each integral by $z_\alpha$ and use the tested probability
law.  The estimates below then have the factor $\prod_\alpha z_\alpha$.
All exceptional contributions may be made smaller than $n^{-A'}$ for an
arbitrary fixed $A'$ by the precision convention in
Definition~\ref{hist:tested-law} and Proposition~\ref{prop:clock-scores}.  If some $z_\alpha<n^{-M}$,
Lemma~\ref{hist:crude-ratios} and the polynomial path moments bound its
unnormalized coefficients by $n^{C(K,B)}z_\alpha$.  Choosing $M$ after $A$
discards all such terms at cost $n^{-A}$.  We shall use this reduction
throughout the proof.

\paragraph{Coefficients with fixed survival and failure weights.}
We start with the terms in which no $w_I$ is differentiated.  Recall from
Proposition~\ref{prop:clock-scores} that, for a node of length $D$,
\begin{equation}\label{repl:score-table}
\begin{aligned}
 |R_l(D)|_\infty&\lesssim_\beta D^{l/2},&
 |R_l(D)|_2&\lesssim_\beta D^{l/2-1},&
 |R_l(D)|_4&\lesssim_\beta D^{l/2-1/2},\\
 |R_3(D)|_1&\lesssim_\beta D^{-1/4},&
 |R_l(D)|_1&\lesssim_\beta D^{(l-4)/2}\quad(l\ge4).
\end{aligned}
\end{equation}
When $D\ge Tn^{-0.005}$, the $L^4$ bound for $R_1$, the $L^2$ bounds for
$R_2,R_3$, and the $L^1$ bound for $R_4$ have additional factors
$n^{-0.01},n^{-0.001},n^{-0.0007}$, respectively.  The notation in
\eqref{repl:score-table} records the length, with the node and its fixed
endpoints understood.

For a last raw coefficient of order $l$ in a forward tested history, the
compensation representations of
Lemma~\ref{score:raw-projection} and Proposition~\ref{prop:clock-scores}(ii), with
$\mathcal J_I\lesssim_\beta1$ by \eqref{score:action-bounds}, give the
following costs on an interval of length
$d$:
\begin{equation}\label{repl:forward-costs}
\begin{array}{c|c}
\text{edge norm}&\text{cost after the compensation identity}\\ \hline
L^2& n^{-1/2}d^{l/2-1/4}\\
L^4& n^{-1/4}d^{l/2-3/16}\\
L^1& n^{-1/2}\max(d^{1/2},d^{l/2-5/4}).
\end{array}
\end{equation}
Each line allows a factor $n^{O(\beta)}$.  These are estimates for the
integrands in the representation after taking the compensator, with earlier
measurable factors kept inside the expectation.  They do not estimate the
expected absolute value of an unprojected score.  The representations used
for different lines or their improvements may be different; the table does
not assert simultaneous bounds on one family of vectors.  The $L^2$ cost gains
$n^{-0.004}$ for $l\ge2$ and $d\ge Tn^{-0.02}$, and the $L^1$ cost at order
four gains $n^{-0.001}$ for $d\ge Tn^{-0.006}$.

Stationary initialization contributes
\[
 \frac{1+u q_e(X_0)}{1+u \bar q_e}.
\]
We need a gain when these are the only differentiated factors.  If $p$ is
a nonnegative stationary initial density with $\mu p=1$ and
$\|p\|_\infty\le n^{C_0}$, then the exact replica identity is
\begin{equation}\label{repl:initial-overlap}
 \frac1{\binom n2}\sum_e \mu[pq_e]^2
 =\frac{n\mu^{\otimes2}[p(X)p(X')\rho_2(X,X')^2]-1}{n-1}
 \le C n^{4\beta}T^{-1}+n^{-A'}
\end{equation}
for every fixed $A'$.  Indeed choose the first dyadic $d_\circ\ge cT$, so
$d_\circ\le2cT$, $r_2(d_\circ)\le c^{-1/2}r_2(T)$, and
$a_2(d_\circ)\ge c_1 n^{3\beta}$.  This scale is among the query scales
in Lemma~\ref{hist:static-support}.  Split at $|\rho_2|=r_2(d_\circ)$;
the complementary probability is
$\exp(-c n^{3\beta})$, which pays $\|p\|_\infty^2$.  This is the weighted
replica identity used in the stationary comparison
\cite[Section~8]{FUA}. %
In particular
\begin{equation}\label{repl:stationary-factor}
 |\bar q|_2\lesssim_\beta T^{-1/2},\qquad
 |\bar q|_4\lesssim_\beta T^{-1/4},\qquad |\bar q|_\infty\le1 .
\end{equation}
The normalized initial marginal of a tested stationary history is such a
$p$: its density is at most $\|v_\alpha\|_\infty/z_\alpha$.  The coefficient
of positive order $j$ in the initialization factor is
$(-\bar q_e)^j+q_e(X_0)(-\bar q_e)^{j-1}$.  A numerator supplies at most one
$q_e(X_0)$ on each history.  Thus if all positive orders come from
initialization, at least one factor is either $\bar q$ or the integrated vector
$\mu[pq_e]$; \eqref{repl:initial-overlap} and boundedness of the other
factors give a fixed power saving.

Suppose now a raw clock likelihood has a positive order.  Each path is a
forest with $O(\log n)$ roots, or consists of the one or two complete legs.
Chronologically earlier roots have lengths at least a constant times later
ones.  Independent paths can be interleaved preserving this property:
merge their lists in decreasing order of the greatest remaining root length,
which is comparable to the current root length on each list.
Project every differentiated root preceding the last such root onto its
endpoints.  The conditional coefficient is
\[
 c_l^I=R_l(I)+q_e(X_{\sup I})R_{l-1}(I),
 \qquad R_0=1,
\]
by \eqref{score:raw-conditional}.  This projection is within the expectation,
using Lemma~\ref{hist:conditional-interior}; the endpoint factors of preceding
roots are measurable before the last root.  After expanding their two
monomials, all initialization coefficients can therefore be bounded by one
in this calculation.

If these preceding roots supply no positive-order $R$ factor, the last
root, of order $l\le k$ and length $d\le CT$, uses the $L^1$ line of
\eqref{repl:forward-costs}.  Its first term is $O(n^{-1/6})$.
Since $T^{3/4}=\sqrt n$, its second term is at most
$CT^{(l-4)/2}$ when $l\ge3$; when $l\le2$ it is at most $T^{-3/4}$,
because $d\ge1$.  For $l<4$ there is a fixed power saving.  Put
\[
 \varepsilon_0=10^{-4}.
\]
For $l=4$ the bound saves $n^{-3\varepsilon_0/4}$ if
$d\le Tn^{-\varepsilon_0}$, while the late $L^1$ improvement saves
$n^{-0.001}$ on the remaining range.  This proves the required bounds for
these terms.

Otherwise, let $p$ be the total order in the preceding positive-order $R$
factors and $D$ their greatest length.  Use the $L^2$ bound on one factor
of length $D$, the supremum bound on the other preceding factors, and the
$L^2$ line of \eqref{repl:forward-costs} for the last order $l$ and length
$d$.  The result is
\begin{equation}\label{repl:unmarked-product}
 n^{O(\beta)}D^{p/2-1}d^{l/2-1/4}T^{-3/4},
 \qquad d\le C D,\qquad p+l\le k.
\end{equation}
For $p+l<4$ this has a fixed power saving.  For $p+l>4$ it is bounded by
$n^{O(\beta)}T^{(p+l-4)/2}$.  To check both assertions when $p=1$, write
the length factor as $(d/D)^{1/2}d^{l/2-3/4}T^{-3/4}$; when $l=1$ it is
even smaller, and otherwise its power of $d$ is positive.

Only $p+l=k=4$ is a borderline case.  Write $D=Tx$, $d=Ty$ with
$y\le Cx$.  Apart from $n^{O(\beta)}$, the three possible factors in
\eqref{repl:unmarked-product} are
\[
 \begin{array}{c|c}
 (p,l)&\text{factor}\\ \hline
 (1,3)&(y/x)^{1/2}y^{3/4}\\
 (2,2)&y^{3/4}\\
 (3,1)&x^{1/2}y^{1/4}.
 \end{array}
\]
If $D$ or $d$ is at most $Tn^{-\varepsilon_0}$, these bounds save at least
$n^{-\varepsilon_0/4}$.  Otherwise every contributing earlier root also
has length at least a constant times $Tn^{-\varepsilon_0}$.  If $l\ge2$,
the late $L^2$ improvement in \eqref{repl:forward-costs} applies.  If
$l=1$ and a preceding order is two or three, choose that factor for the
$L^2$ bound and use its $n^{-0.001}$ improvement.  Using its length in
place of the greatest length costs at most $n^{\varepsilon_0/2}$, since
all those lengths lie between a constant times $Tn^{-\varepsilon_0}$
and $CT$.  Finally, if the preceding orders are $(1,1,1)$, apply
$L^4$ to all four factors.  The three $R_1$ factors have their late
improvements, and the last factor is at most
$n^{O(\beta)}T^{5/16}n^{-1/4}=n^{-1/24+O(\beta)}$.
These are fixed power savings in all cases.  The $O(\log n)^K$ choices of
roots are absorbed into $n^{O(\beta)}$.  This proves
\eqref{repl:averaged-history} for terms with fixed cutoff weights.

\paragraph{The exact expansion at a differentiated cutoff.}
For a node $I=[a,b]$, use the pre-density $\bar k_I$ and conditional
interior law $\mathsf C_I^{x,y}$ from Lemma~\ref{hist:conditional-interior}.
For the perturbation at $e$, put
\[
 H_{I,e}(u)=\frac{\bar k_{I,e}(u)}{\bar k_I},\quad H_{I,e}(0)=1,
 \qquad
 V_{I,e}^+(u)=\frac{w_{I,e}(u)}{w_I},\quad
 V_{I,e}^-(u)=\frac{1-w_{I,e}(u)}{1-w_I}.
\]
All denominators are positive.  Write $W_N^+=W_N$ and $W_N^-=1-W_N$.
For $j\ge1$, the scalar chain rule gives
\begin{equation}\label{repl:cutoff-chain}
\begin{split}
 [u^j]V_{I,e}^{\sigma}(u)
 &=\sum_{\ell=1}^j
 \frac{(W_{N_I}^{\sigma})^{(\ell)}(\log\bar k_I)}
      {\ell!\,W_{N_I}^{\sigma}(\log\bar k_I)}
 \sum_{\substack{a_1+\cdots+a_\ell=j\\a_h\ge1}}
       \prod_{h=1}^{\ell} b_{I,e,a_h},\\
 b_{I,e,a}
 &:=[u^a]\log H_{I,e}(u)\\
 &=\sum_{v=1}^a\frac{(-1)^{v+1}}v
       \sum_{\substack{d_1+\cdots+d_v=a\\d_h\ge1}}
       \prod_{h=1}^{v}[u^{d_h}]H_{I,e}(u),
 \qquad \sigma\in\{+,-\}.
\end{split}
\end{equation}
Thus a positive-order cutoff coefficient supplies one positive-order
derivative ratio of $W_N^\sigma$ and a product of positive-order pre-density
coefficients whose orders sum to its assigned order.  We call this occurrence
of a derivative ratio a \emph{cutoff mark} at $I$.

The pre-density ratio has the exact representation
\begin{equation}\label{repl:pre-density-copy}
 H_{I,e}(u)
 =(1+u \bar q_e)\,
 \mathsf C_I^{x,y}\left[
    \frac{L_{I,e}(u)}{1+u q_e(y)}
    \prod_{\substack{J\subsetneq I\\J\text{ a checked node}}}
           V_{J,e}^+(u)\right].
\end{equation}
Indeed, the perturbed stationary probability is
$\mu_u(y)=\mu(y)(1+u q_e(y))/(1+u \bar q_e)$.  Dividing the perturbed
pre-truncation transition mass by its base value gives the conditional
expectation of $L_{I,e}(u)\prod_{J\subsetneq I}V_{J,e}^+(u)$.
The further factor $\mu(y)/\mu_u(y)$ gives
\eqref{repl:pre-density-copy}; see also Equation~\eqref{hist:pre-ratio}.
It is essential here that $\bar k_I$ is a density with respect to the
normalized stationary probability.

Each product of coefficients of $H_{I,e}$ in \eqref{repl:cutoff-chain}
is represented using independent interiors with law $\mathsf C_I^{x,y}$
conditional on the same endpoints.  Recursively do the same for every
positive-order descendant cutoff coefficient in
\eqref{repl:pre-density-copy}.  Given all already sampled paths and the
endpoints of a new copy, its interior is sampled independently with this
law.  Ancestor and exterior factors depend only on data outside the copied
interior.  By Lemma~\ref{hist:conditional-interior}, replacing that interior in
the full tested history preserves its law.  This remains true for a copy
inside a copy by iteration.  In particular the marginal estimates for each
copy are those for a base tested path, averaged over its endpoints; no
pointwise assertion for arbitrary endpoints is being made.

There is a useful order invariant in this expansion.  Assign to each
coefficient slot its positive Taylor order.  Formula
\eqref{repl:cutoff-chain} splits an order into positive orders with the
same sum, and expanding a slot by \eqref{repl:pre-density-copy} has the
same property.  The total order in all unexpanded slots, raw likelihood
slots, and endpoint or $\bar q_e$ factors is therefore at most $k$ at every
stage.  Cutoff marks record transferred orders; they do not add to this sum.
There may be a long chain of cutoff marks carrying a single order.

A second invariant gives a projected factor when expansion terminates.
For a positive-order pre-density slot of order $r$ whose descendants receive
no positive cutoff order, keep the endpoint denominator together with the
conditional likelihood.  Its coefficient is exactly
\begin{equation}\label{repl:terminal-mark}
 [u^r](1+u \bar q_e)R_e(I;u)
       =R_r(I)+\bar q_e R_{r-1}(I).
\end{equation}
Each summand contains a positive-order $R_l(I)$, or the order-one factor
$\bar q_e$ (when $r=1$ the second summand is $\bar q_e$ itself).
Induction from terminal slots shows that every terminal branch descending
from a cutoff mark contributes a positive-order conditional coefficient or an
$\bar q_e$ factor.  The endpoint denominator in
\eqref{repl:terminal-mark} is never estimated separately.  By
\eqref{repl:stationary-factor}, $\bar q_e$ can be treated as an $R_1$ factor
of length $T$: it satisfies all the corresponding bounds in
\eqref{repl:score-table}, including the late $L^4$ improvement.

\paragraph{Summation over cutoff marks and termination of the expansion.}
Let $\lambda_I=1-w_I$ for a survival cutoff mark and $\lambda_I=1$ for a
failure cutoff mark.  For any fixed $\eta>0$,
Lemma~\ref{hist:cutoff-hazards} gives, outside an arbitrarily small
inverse-polynomial contribution,
\begin{equation}\label{repl:mark-gain}
 \left|
 \frac{(W_{N_I}^{\sigma})^{(\ell)}(\log\bar k_I)}
      {\ell!\,W_{N_I}^{\sigma}(\log\bar k_I)}
 \right|
 \le C n^\eta\frac{\lambda_I}{N_I},
 \qquad
 N_I=n^\beta(T/|I|)^{3/2}.
\end{equation}
For survival this estimate is used on $w_I\ge n^{-L}$.
The discarded event $w_I<n^{-L}$ is paid from its included base survival
factor, with $L$ chosen arbitrarily large.  For failure the estimate holds
without this restriction.  The global derivative-ratio bounds of
Lemma~\ref{hist:cutoff-hazards} control exceptional contributions.

Here is a precise way to sum the choices of nodes in
\eqref{repl:mark-gain}.  On a normalized tested history with mass
$z_\alpha\ge n^{-M}$ and test bounded by $n^B$, put
$H=\sum_{I\in\mathscr S_\alpha}(1-w_I)$.  Since
$\prod_I w_I\le e^{-H}$,
\[
 \mathbb P_{\mathrm{test}}\{H>h\}\le n^{M+B}e^{-h}.
\]
One possible failure adds at most one to this hazard sum.  The same
statement holds for the marginal of every conditional copy, by its
law-preserving construction.  To make these statements simultaneous, for
each deterministic prefix of at most $P+K$ choices of nodes carrying cutoff
marks, sample the possible next copies conditional on that prefix.  All continuations
of a prefix share its already sampled paths.  There are $n^{O_K(P)}$
prefixes for fixed $P$.  The arbitrary precision convention permits the
hazard and moment bounds to hold for all of them simultaneously.  This union
uses the joint copy law and the marginal estimates for each indexed path.

All expectation and compensator identities used for this summation are
established in their stated filtrations before restriction to the
simultaneous event.  At a fixed sampled prefix, let $\mathcal I\ne\varnothing$
be the eligible next nodes on one shared path.  For each $I$, let $B_I\ge0$
bound the absolute value of the resulting continuation integrand or its edge norm on the common
sampled tree, including all later factors and gains but excluding the factor
$C n^\eta\lambda_I/N_I$ used to bound the current mark ratio.  The bound $B_I$
may depend on later sampled data.  On this common realization,
\[
 \sum_{I\in\mathcal I}\frac{\lambda_I}{N_I}B_I
 \le \left(\sum_{I\in\mathcal I}\lambda_I\right)
          \sup_{I\in\mathcal I}\frac{B_I}{N_I}.
\]
On the simultaneous event, the parenthesis is at most $n^\eta+1$.
Combining this with the current $C n^\eta$ from \eqref{repl:mark-gain}
costs at most $C n^{2\eta}$ while retaining $N_I^{-1}$ with the continuation
for the same choice.  Apply this bound separately to eligible slots and paths,
paying the fixed combinatorial multiplicity.  Iteration gives a supremum over
indexed full branches.  On the simultaneous event, the score bounds together
with the length estimates below control this supremum before integration.  The hazards are
scalars independent of $e$, so the same argument applies to edge-norm
estimates after projection.

The indicator of the simultaneous event is used only to bound these
resulting integrands.  The complementary integrals are
controlled by the polynomial moments in Lemmas~\ref{hist:crude-ratios}
and \ref{score:hellinger} and in the proof of
Lemma~\ref{score:raw-projection}, choosing their precision after $P$.
Thus no future-dependent hazard event is used
as a predictable factor.

It suffices to stop after a fixed number $P$ of cutoff marks.  If this
threshold is reached, leave the remaining pre-density slots unexpanded.  By
Lemma~\ref{hist:crude-ratios} and the fixed-order path moment bounds, their
product, together with all other positive-order slots, is at most
$n^{C_{\mathrm{cr}}}$ in the normalized estimate, for a constant
$C_{\mathrm{cr}}=C_{\mathrm{cr}}(K)$ independent of $P,\beta$.
Indeed the sum of their positive orders is at most $K$; the derivative
ratios already extracted at cutoff marks have uniform global bounds.  Every
node has length at most $CT$, so $N_I^{-1}\le Cn^{-\beta}$.  The contribution
stopped at $P$ cutoff marks, with an overshoot of at most $K$, is therefore
at most
\begin{equation}\label{repl:depth-remainder}
 C_{K,P}n^{C_{\mathrm{cr}}-P\beta+2(P+K)\eta}
\end{equation}
apart from exceptional precision.

The parameter choices can now be made in order.  First $K=24$ is fixed.
After choosing a sufficiently small $\beta$ for the score estimates, take
$P>(C_{\mathrm{cr}}+3)/\beta$, and then
$\eta\le\beta/(4(P+K))$.  For this fixed $P$, all scalar multiplicities
from \eqref{repl:cutoff-chain} and powers of $\log n$ from partitions cost
at most $n^{\beta/2}$ for large $n$.  Together with
$2(P+K)\eta\le\beta/2$, \eqref{repl:depth-remainder} is at most $n^{-2}$.
Finally choose the exceptional-set and small-mass precisions.  In a fully
expanded term at most $K$ positive-order factors use the score estimates.
Their total exponent loss is thus $C_K\beta$ independent of $P$.

\paragraph{The edge norms in fully expanded terms with cutoff marks.}
It remains to estimate terms which terminate before depth $P$.  We first
form a partition from the locations of their cutoff marks, before deciding
which endpoints to condition on.  Split each root into dyadic pieces only
at nodes which properly contain a descendant carrying a cutoff mark.  Make
the same partition on a copy using its own cutoff marks.

Let $m$ be the least length of a node carrying a cutoff mark, capped at $T$.
In each fully expanded monomial, select one factor $R_l(I)$ with $l\ge1$, or
one factor $\bar q_e$, supplied by a terminal $H_{I,e}$ slot through
\eqref{repl:terminal-mark}.  Make this choice from the indexed monomial,
hence the same way for every $e$.  The slot was created by a cutoff mark at
$I$, so the selected factor has length $|I|\ge m$ in the $R_l(I)$ case and
assigned length $T\ge m$ in the $\bar q_e$ case.  Every other terminal factor
remains in the product.  The gain from a shortest cutoff mark will be retained
separately below.  All partition pieces have length at
least a constant times $m$, except possibly untouched original forest
roots.  Indeed, a split made to isolate a descendant $M$ carrying a cutoff
mark creates two children, each of length at least $|M|\ge m$; further
splits satisfy the same statement with their own such descendants.  A copy
starts at a node carrying a cutoff mark, so it has no untouched smaller
root.  This also gives $O_P(1+\log n)$ pieces.  The untouched roots smaller
than a sufficiently small constant times $m$ can occur only in a stopped
forest, and occur later than all cutoff marks on that path: the lengths of
its original roots decrease chronologically.  Choose that constant below
the constants for complete legs, so no complete leg is included in this
exception.

Projecting the raw order-one coefficient on a root of length $D_1$ would
replace it by $c_1=R_1(D_1)+q_{\mathrm{end}}$.  If paired with $R_3(D_3)$,
the $R_3(D_3)R_1(D_1)$ summand would have $L^2$-times-$L^2$ cost
$n^{O(\beta)}(D_3/D_1)^{1/2}$ by \eqref{repl:score-table}, which can grow
when $D_1$ is very short.  We therefore reserve one allocation before
projecting piece coefficients.
An untouched original root $I$ of length $z<cm$ may carry a raw coefficient
of order one while the sole other factor of positive order is a conditional
coefficient of order three.
This is the small-root $(3,1)$ allocation.  Its terminal state
is not revealed in advance, and its complete raw coefficient
$\ell_{1,e}(I)=[u]L_{I,e}(u)$ is retained until forward compensation.
The total order is four, so no other positive-order factor is present.
The small root carries no cutoff mark and cannot supply the selected terminal
factor; hence in this allocation the selected factor is the $R_3$ factor.

For every other allocation, condition on the relevant endpoints of nodes
carrying cutoff marks and on the partition endpoints.  The fresh piece
interiors then have their conditional laws $\mathsf C_J$ and factor after
the endpoint variables and their shared identifications are fixed: lower
weights stay inside their pieces, and every other included weight depends
only on the conditioned endpoints.  Shared exterior variables occur once;
this does not assert independence of the completed copied histories.
Whenever a raw coefficient on a piece is projected, it becomes
$c_l=R_l+q_{\mathrm{end}}R_{l-1}$.  Nonterminal denominators in
\eqref{repl:pre-density-copy} and initialization may contribute endpoint
powers of absolute value one.  At a terminal slot use
\eqref{repl:terminal-mark} instead.  The resulting terms are products of
$R$ factors, $\bar q$ factors counted as $R_1(T)$, and bounded endpoint
factors.  The sum of their positive orders is at most $k$, and the selected
terminal factor remains in this product.  All these conditional identities
are integrated under the joint copy law when the sharp estimates are used;
they give no uniform sharp bound at an arbitrary fixed endpoint graph.

For total $R$ order at most four, \eqref{repl:score-table} gives the
following complete list.  Unused H\"older exponents are filled with the
constant function on the edge probability space.
\begin{equation}\label{repl:low-patterns}
\begin{array}{c|c|c}
\text{orders of the }R\text{ factors}&\text{norms used}
   &\text{bound apart from }n^{O(\beta)}\\ \hline
1^r,\ 1\le r\le4& L^4\text{ for each}&1\\
2,1^r,\ 0\le r\le2&L^2,L^4,\ldots,L^4&1\\
2,2&L^2,L^2&1\\
3&L^1&D_3^{-1/4}\\
4&L^1&1\\
3,1&L^2,L^2&(D_3/D_1)^{1/2}.
\end{array}
\end{equation}
Thus only $(3,1)$ can cost more than a constant.  If its order-one length
$D_1$ is at least a constant times $m$, its cost is at most
$C(T/m)^{1/2}$.  Retain the gain from a shortest cutoff mark,
\begin{equation}\label{repl:minimum-mark-gain}
 \frac1{N(m)}=n^{-\beta}(m/T)^{3/2};
\end{equation}
when the shortest node carrying a cutoff mark has length greater than $T$,
the same formula holds up to a constant.  The remaining gains are at most a constant times
$n^{-\beta}$ each and pay their hazard summations.
If $m\le Tn^{-\varepsilon_0}$,
\eqref{repl:minimum-mark-gain} leaves at most $(m/T)$ in the $(3,1)$
case and $(m/T)^{3/2}$ in all other cases, after the total $n^\beta$
allowance for hazards and logarithms.  This is a fixed power saving.

For the reserved small-root allocation, retaining $\ell_{1,e}(I)$ keeps its
$R_1$ and $q_{\mathrm{end}}$ conditional terms together before a triangle
inequality.  The selected order-three factor is earlier on the same path or
belongs to an independent original history.  The endpoints of all nodes carrying
cutoff marks on the same path precede $I$.  Their conditional copies can be
sampled before $I$, independently of its future given those already observed
endpoints.  Factors on other original histories can also be exposed first,
because those histories and their external tests factor separately.  This
auxiliary sampling is fixed as part of the admissible law before choosing
the earlier vector factor.

After projecting the earlier pieces, we return to the joint tested
expectation with the terminal state of $I$ unrevealed in advance, and
compensate its raw order-one coefficient in the forward filtration.
Under the forward tested law on this root, the external terminal and
possible failure tests are included in its tilted intensities.
The $L^2$ compensation in \eqref{repl:forward-costs} is therefore valid
with these earlier factors retained and gives
$n^{O(\beta)}z^{1/4}/\sqrt n$.  Paired with the order-three $L^2$ bound,
the result is
\begin{equation}\label{repl:small-root}
 n^{O(\beta)}\frac{D_3^{1/2}z^{1/4}}{\sqrt n}
 =n^{O(\beta)}(D_3/T)^{1/2}(z/T)^{1/4}
 \le n^{O(\beta)}.
\end{equation}
This calculation is made before projecting the small root onto its
endpoint, and the compensation identity precedes restriction to the
hazard event.  It proves the same early-$m$ saving using
\eqref{repl:minimum-mark-gain}.

For completeness, suppose $m\ge Tn^{-\varepsilon_0}$.  The selected terminal
factor has length at least $m$ and is in the late range.  Except for $(3,1)$,
the norms in \eqref{repl:low-patterns} give a fixed improvement at that
factor: a selected $R_1$ uses its late $L^4$ bound, a selected $R_2$ its late
$L^2$ bound, and a selected $R_4$ its late $L^1$ bound.
If the selected factor is the singleton $R_3$, then
$D_3^{-1/4}\le T^{-1/4}n^{\varepsilon_0/4}$.
For the regular $(3,1)$ case, put $\delta_0=0.005$ and use the
$R_3$ late $L^2$ gain $n^{-0.001}$.  Since $D_1\ge cm$,
\[
 \begin{cases}
 n^{-0.001}(D_3/D_1)^{1/2}
       \le C n^{-0.001+\varepsilon_0/2},
       &D_3\ge Tn^{-\delta_0},\\
 (D_3/D_1)^{1/2}
       \le C n^{-(\delta_0-\varepsilon_0)/2},
       &D_3<Tn^{-\delta_0}.
 \end{cases}
\]
The two savings are $n^{-0.00095}$ and $n^{-0.00245}$.
In \eqref{repl:small-root}, the corresponding cases gain $n^{-0.001}$
and $n^{-\delta_0/2}$, respectively.  These margins, and the early
$m$ margin, remain positive after an exponent $C_K\beta$ is included.
Together with the fixed-weight calculation, they prove the low-order
part of \eqref{repl:averaged-history}; one may take, for example,
$a=10^{-7}$.

For $k>4$ the product estimates need no late gains or lower length bound.
The following allocation of H\"older norms makes the power count explicit.
A factor of order $l\ge4$ uses its $L^1$ bound and saves two powers of
$T$ relative to the supremum bounds with all lengths replaced by $CT$.
If the largest order
is three, pair it with an order two or three in $L^2$, or with two order
ones using $L^2,L^4,L^4$; again the saving is two powers of $T$.
A singleton order three is bounded by one.  An order three paired with
exactly one order one costs at most $T^{1/2}$, which is allowed because
$k\ge5$.  If all orders are at most two, two order-two factors use
$L^2,L^2$, one order two and two order ones use $L^2,L^4,L^4$, and four
order ones use $L^4$ four times.  Each of these saves two powers of $T$;
when fewer factors are present their product is bounded by one.
All unused factors use their supremum bounds.  Thus if the total $R$
order is $h\le k$, the cost is at most
$n^{O(\beta)}T^{(k-4)/2}$.  This proves the high-order part.

Finally the same expansion with the supremum bounds in
Proposition~\ref{prop:clock-scores}, the raw fixed-order path bounds, and
$\sum l\le k$ gives \eqref{repl:sup-history}.  The depth remainder is
already smaller, and the exceptional and small-mass reductions at the
start give the stated additive precision.  Since the Taylor coefficients
of $\tanh\vartheta$ at zero are fixed constants, a derivative of order $k$
in $\vartheta$ is a fixed linear combination of coefficients of orders at
most $k$ in $u$.  Applying this to each factor of a product leaves total
$u$ order at most the assigned total $\vartheta$ order.  The lower-order
estimates fit the asserted bounds for that total order.
\end{proof}

\subsection{Localized traces and entrywise comparison}

Let $\mathcal Q_n$ be a finite set of queries in the fixed compact
rectangles, with $|\mathcal Q_n|\le n^{d_0}$ for a fixed $d_0$.  It may
contain both stationary and quench queries.  For a query $q$, let
$\zeta_q$ be its endpoint observable
$n^{-2/3}\sum_i X_iY_i$, so $|\zeta_q|\le n^{1/3}$.  Define
\[
 \widetilde Y_q=
 \mathbb E_q\left[\zeta_q\prod_{I\text{ in }q}w_I\right],
 \qquad
 \ell_q=1-\mathbb E_q\left[\prod_{I\text{ in }q}w_I\right].
\]
The expectations have the actual initialization of the query.  In
particular the quench history starts from the uniform cube law.
The surviving trace $\widetilde Y_q$ is not renormalized.
Order checks chronologically as in Definition~\ref{hist:tested-law}.  Telescoping
the product gives the exact first-failure decomposition
\begin{equation}\label{repl:first-failure}
 \ell_q=\sum_{I\text{ in }q}\ell_{q,I},\qquad
 \ell_{q,I}:=\mathbb E_q\left[
       \prod_{J\text{ preceding }I}w_J\,(1-w_I)\right]\ge0 .
\end{equation}
The summand is the mass of the corresponding checked first-failure prefix;
its unneeded future integrates to one.

Let $\Gamma(J)=\chi_n((J_e/\sqrt n)_e)$ be the static cutoff of
Lemma~\ref{hist:static-support}.
Choose a smooth $\chi:\mathbb R\to[0,1]$ equal to one on $(-\infty,1]$
and zero on $[2,\infty)$, and put
\begin{equation}\label{repl:outer-cutoff}
 G(J)=\Gamma(J)\chi\left(n^{10}\sum_{q\in\mathcal Q_n}\ell_q(J)\right).
\end{equation}
The Gaussian static estimate and Lemma~\ref{hist:gaussian-survival}, used at a
precision chosen after $d_0$, imply
\begin{equation}\label{repl:gaussian-cutoff}
 \mathbb E_{\mathrm{G}}G\longrightarrow1.
\end{equation}

We compare bounded smooth functions $F=F_n$ on
$\mathbb R^{\mathcal Q_n}$ satisfying, for fixed constants independent of
$n$,
\begin{equation}\label{repl:test-class}
 \|F\|_\infty\le C_F,\qquad
 \sup_y\sum_{q_1,\ldots,q_j\in\mathcal Q_n}
 \left|\partial_{q_1}\cdots\partial_{q_j}F(y)\right|
 \le C_F(\log n)^{C_F},\quad 1\le j\le K .
\end{equation}
This class includes $F=1$.  Write
$\Phi(J)=G(J)F((\widetilde Y_q(J))_{q\in\mathcal Q_n})$ and define
\[
 \partial_e^j\Phi(J)
 =\left.\frac{d^j}{d\vartheta^j}
       \Phi(J_{-e},J_e+\sqrt n\,\vartheta)\right|_{\vartheta=0}.
\]
Thus $\partial_e$ differentiates the interaction coefficient, rather than
the unscaled entry $J_e$.

\begin{lemma}[Derivatives of the localized observable]\label{repl:observable-derivatives}
The parameter $\beta>0$ can be chosen so that there are
$a_0>0$ and $0<\delta<0.01$ with, for every disorder matrix $J$,
\begin{align}
 \frac1{\binom n2}\sum_e|\partial_e^4\Phi(J)|
     &\le n^{-a_0},                                           \label{repl:fourth-derivative}\\
 \frac1{\binom n2}\sum_e|\partial_e^k\Phi(J)|
     &\le n^\delta T^{(k-4)/2},\quad 5\le k<K,               \label{repl:higher-derivative}\\
 \max_e|\partial_e^j\Phi(J)|
     &\le n^\delta T^{j/2},\quad 1\le j\le K.                \label{repl:global-derivative}
\end{align}
The constants may depend on the fixed exponents in
\eqref{repl:test-class} and on $d_0,c,C$.
\end{lemma}

\begin{proof}
On the support of a term containing a differentiated history, the static
support conditions hold and
$\sum_q\ell_q\le2n^{-10}$.  This follows from the supports of $\Gamma$
and its derivatives, and of $\chi$ and its derivatives.  Split each
$\zeta_q$ into its positive and negative parts.  The base masses of these two
checked integrals are at most $n^{O(\beta)}$.  Indeed the post-check kernel
caps at the ends of the one or two legs dominate the joint law of the two observation spins
by
$\exp(C'N(cT))\mu^{\otimes2}$.  For a stationary query the first spin
already has law $\mu$; for a uniform-start query the first leg cap holds
from each initial spin.  Choose the first dyadic $d_*\ge cT$.  Then
$d_*\le2cT$, $r_2(d_*)\le r_2(cT)$, and $a_2(d_*)\ge c_1n^{3\beta}$.
The pair cutoff therefore gives
\[
 \mathbb E_q\left[|\zeta_q|\prod_Iw_I\right]
 \le n^{1/3}r_2(cT)
     +n^{1/3}\exp(C'N(cT)-c n^{3\beta})
 \le C n^{2\beta}+n^{-A'}
\]
for arbitrary fixed $A'$.  We used $N(cT)=O(n^\beta)$.

Apply the ordinary chain rule to $\chi(n^{10}\sum\ell_q)$ and to
$F(\widetilde Y)$.  A term has at most $K$ positive-order history
coefficients, with total order at most the differentiating order; repeated
histories are represented by independent copies.  The coefficients from
derivatives of $F$ are evaluated at the base $\widetilde Y(J)$, hence are
common across $e$, and their total absolute sums are bounded by
\eqref{repl:test-class}.  For loss factors, expand
\eqref{repl:first-failure}.  The mass factors in
Lemma~\ref{repl:history-derivatives}, summed over all choices of queries and
first-failure nodes and multiplied by their $n^{10}$ factors, give at most
\[
 \left(n^{10}\sum_{q,I}\ell_{q,I}\right)^r\le2^r
\]
when there are $r$ loss factors.  Trace mass factors cost at most
$n^{O_K(\beta)}$.  The additive precisions in that lemma are chosen after
the polynomial number of nodes and queries and after the factors
$n^{10K}$, so their total is negligible.

If a static cutoff is also differentiated, its pointwise derivative bound
from Lemma~\ref{hist:static-support} multiplies these estimates.  If only the
static cutoff is differentiated, that bound is $O(n^{-b})$ for some
$b=b(\beta)>0$, and its remaining outer factor is bounded.  Consequently
the fourth averaged derivative is at most
$n^{-a+C'_K\beta}(\log n)^{C'}+O(n^{-b})$, higher averaged derivatives
are at most $n^{C'_K\beta}(\log n)^{C'}T^{(k-4)/2}$, and the supremum
bound is at most $n^{C'_K\beta}(\log n)^{C'}T^{j/2}$.
These estimates hold everywhere: outside the indicated supports the
corresponding product-rule terms vanish.  Choose $\beta$ after $K$ so
that $C'_K\beta<a/2$ and $C'_K\beta<0.005$, while satisfying the preceding
sections' smallness requirements.  Absorb logarithms by a still smaller
fixed power.  This gives \eqref{repl:fourth-derivative}--%
\eqref{repl:global-derivative} for some $a_0>0$ and $\delta<0.01$.
\end{proof}

\begin{lemma}[Entrywise replacement]\label{repl:comparison}
For every sequence of query sets and tests satisfying
\eqref{repl:test-class},
\[
 \mathbb E_{\mathrm{Rad}}\!\left[G F(\widetilde Y)\right]
 -\mathbb E_{\mathrm{G}}\!\left[G F(\widetilde Y)\right]\longrightarrow0 .
\]
\end{lemma}

\begin{proof}
We use the Taylor replacement principle \cite{Chatterjee2006}, in the finite
mixture form of the stationary comparison \cite[Section~8]{FUA}.
Put
$\nu_\lambda=(1-\lambda)\mathsf N(0,1)
 +\lambda(\delta_{-1}+\delta_1)/2$ for $0\le\lambda\le1$.
Differentiating its finite product measure gives exactly
\begin{equation}\label{repl:mixture-derivative}
 \frac d{d\lambda}\mathbb E_{\nu_\lambda^{\otimes\binom n2}}\Phi
 =\sum_e\mathbb E_{\lambda,-e}
    \left[\mathbb E_\xi\Phi(J_{-e},\xi)
                 -\mathbb E_Z\Phi(J_{-e},Z)\right],
\end{equation}
where $\xi$ is a symmetric sign, $Z$ is standard normal, and
$\mathbb E_{\lambda,-e}$ averages the other entries.

Independently of $J_{-e},\xi,Z$, sample $V$ from $\nu_\lambda$ and set
$J^0=(J_{-e},V)$.  Taylor-expand at this completed matrix:
\begin{equation}\label{repl:centered-taylor}
 \Phi(J_{-e},x)=
 \sum_{k=0}^{K-1}
   \frac{\partial_e^k\Phi(J^0)}{k!\,n^{k/2}}(x-V)^k
       +\mathcal R_e(x,V),
 \qquad
 |\mathcal R_e(x,V)|
 \le\frac{|x-V|^K}{K!\,n^{K/2}}\,n^\delta T^{K/2}.
\end{equation}
The remainder bound is global by \eqref{repl:global-derivative}, including
every point of the segment.  Conditional on $J_{-e},V$, the orders
$k\le3$ cancel exactly between the two proposals.  In fact
\begin{equation}\label{repl:centered-moments}
 \Delta_k(V):=\mathbb E_\xi(\xi-V)^k-\mathbb E_Z(Z-V)^k
 =\sum_{j=4}^k\binom kj(-V)^{k-j}
      \bigl(\mathbb E\xi^j-\mathbb EZ^j\bigr),
\end{equation}
where the sum is empty for $k\le3$.  Independence of the completed entry
from both proposals is what permits this conditional cancellation.

For $4\le k<K$, on $|V|\le\log n$ the absolute value of
\eqref{repl:centered-moments} is at most $C_K(\log n)^K$.  The complementary
mixture tail, even with any fixed moment of $V$, is smaller than every
inverse power of $n$.  Its contribution is negligible by the global
polynomial bounds \eqref{repl:global-derivative}.  When $e$ is uniform,
$J^0$ has the full iid mixture law independently of $e$.  Therefore
\eqref{repl:mixture-derivative}--\eqref{repl:centered-taylor} yield,
uniformly in $\lambda$,
\begin{align}
 \left|\frac d{d\lambda}\mathbb E_{\nu_\lambda^{\otimes\binom n2}}\Phi\right|
 &\le C_K(\log n)^K
   \sum_{k=4}^{K-1}n^{2-k/2}
     \mathbb E_\lambda\!\left[
       \frac1{\binom n2}\sum_e|\partial_e^k\Phi(J)|\right]\notag\\
 &\hspace{2cm}+C_K n^{2-K/2+\delta}T^{K/2}+o(1).             \label{repl:master-comparison}
\end{align}
The remainder uses the bounded $K$th moments of $\xi,Z,V$.
By \eqref{repl:fourth-derivative}, the $k=4$ term tends to zero.
For $k\ge5$, \eqref{repl:higher-derivative} and $T=n^{2/3}$ give
\[
 n^{2-k/2+\delta}T^{(k-4)/2}
      =n^{-(k-4)/6+\delta}=o(1).
\]
The remainder is $n^{2-K/6+\delta}=n^{-2+\delta}$ because $K=24$.
Thus \eqref{repl:master-comparison} tends to zero uniformly in $\lambda$.
Integration over $[0,1]$ proves the lemma.
\end{proof}

\subsection{Removing the cutoff and transferring tightness}

Take $F=1$ in Lemma~\ref{repl:comparison}.  In view of
\eqref{repl:gaussian-cutoff},
\begin{equation}\label{repl:both-cutoffs}
 \mathbb E_{\mathrm{Rad}}(1-G)\longrightarrow0,
 \qquad
 \mathbb E_{\mathrm{G}}(1-G)\longrightarrow0.
\end{equation}
This step transfers the high probability of the cutoff to Rademacher
disorder using derivative estimates that held for every matrix on their
deterministic supports.

Let $Y_q$ be the actual trace.  On $G>0$,
\begin{equation}\label{repl:remove-loss}
 |Y_q-\widetilde Y_q|
 \le n^{1/3}\ell_q
 \le2n^{1/3-10},\qquad q\in\mathcal Q_n.
\end{equation}
The first-derivative bound in \eqref{repl:test-class} makes $F$ Lipschitz
for the coordinate supremum norm with constant $C_F(\log n)^{C_F}$.
It follows from \eqref{repl:both-cutoffs}--\eqref{repl:remove-loss} that,
for either entry law,
\[
 \mathbb E F(Y)-\mathbb E[G F(\widetilde Y)]\longrightarrow0.
\]
Hence Lemma~\ref{repl:comparison} applies to the actual traces at finite query
sets and at all polynomial grids with tests in
\eqref{repl:test-class}.  At fixed finite sets this proves the equality
of limiting joint finite-dimensional laws with the Gaussian pair.

We give the tightness transfer because a comparison only at fixed times
would not suffice for the asserted topology.  The finite heat-bath
generator $\mathcal L_J=\sum_i(P_i-I)$ has
$\|\mathcal L_J\|_{\infty\to\infty}\le2n$, and its semigroup is a
contraction in supremum norm.  Differentiating the one or two semigroups
in a trace therefore gives, for every $J$ and on the fixed rectangles,
\begin{equation}\label{repl:time-lipschitz}
 |\partial_t A_{n,J}(t)|\le2n^2,\qquad
 |\partial_s B_{n,J}(s,t)|+|\partial_t B_{n,J}(s,t)|\le4n^2.
\end{equation}
Indeed each derivative in $s$ or $t$ inserts $T\mathcal L_J$ into its
corresponding semigroup, and the undifferentiated trace is bounded by $n^{1/3}$.

Choose grids of mesh at most $h_n=n^{-4}$ on the two rectangles.  They and
all pairs of their points have polynomial size, and
\eqref{repl:time-lipschitz} makes their interpolation error $O(n^{-2})$.
For grid coordinates $y$, the signed forms used to test a supremum are
$\pm y_q$, and those for a modulus of continuity are
$\pm(y_q-y_{q'})$ for pairs in the same rectangle at distance at most
a specified radius.  Include the zero form.  Each such linear form has
coefficient $\ell^1$ norm at most two.  For a collection of $M_n$ forms
$a_r(y)$, with $M_n$ polynomial, put
\[
 \operatorname{smax}_{\lambda_n}(a(y))
 =\lambda_n^{-1}\log\sum_{r=1}^{M_n}e^{\lambda_n a_r(y)},
 \qquad \lambda_n=(\log n)^2.
\]
It lies between $\max_r a_r$ and
$\max_r a_r+(\log M_n)/\lambda_n$, so its error is $o(1)$.
Its derivatives obey
\begin{equation}\label{repl:softmax-derivatives}
 \sum_{q_1,\ldots,q_j}
 \left|\partial_{q_1}\cdots\partial_{q_j}
      \operatorname{smax}_{\lambda_n}(a(y))\right|
 \le C_j\lambda_n^{j-1},\qquad 1\le j\le K.
\end{equation}
For example, the derivative of order $j$ is $\lambda_n^{j-1}$ times
the joint cumulant of the coefficient coordinates of a random form chosen
with its normalized exponential weight.  The cumulant partition formula,
followed by the coefficient $\ell^1$ bound two in each factor, proves
\eqref{repl:softmax-derivatives}.  Composing with a fixed bounded smooth
function whose derivatives through order $K$ are bounded now satisfies
\eqref{repl:test-class} by the ordinary chain rule.

Let $\omega_n(\rho)$ be the greater of the two moduli of continuity on
the fixed rectangles, with the supremum metric on each parameter domain.
Use the preceding smoothed maximum over grid pairs at distance at most
$2\rho$, followed by a smooth increasing function which is zero below
$\epsilon/3$ and one above $2\epsilon/3$.  For fixed $\rho,\epsilon>0$,
the interpolation and smoothing errors are smaller than $\epsilon/12$
for large $n$.  Comparison of the actual grid tests consequently gives
\begin{equation}\label{repl:modulus-transfer}
 \limsup_{n\to\infty}\mathbb P_{\mathrm{Rad}}\{\omega_n(\rho)>\epsilon\}
 \le
 \limsup_{n\to\infty}\mathbb P_{\mathrm{G}}\{\omega_n(2\rho)>\epsilon/4\}.
\end{equation}
Here a pair of original points at distance at most $\rho$ is moved to
nearest grid points at distance at most $2\rho$.  The same construction
with forms $\pm y_q$ transfers tightness of the supremum norm.  Gaussian
tightness, proved in Proposition~\ref{prop:gaussian-subsequence} and Theorem~\ref{thm:entrance-selection},
makes the right side of \eqref{repl:modulus-transfer} tend to zero as
$\rho\downarrow0$ and makes the supremum tails vanish as their threshold
tends to infinity.  The Arzel\`a--Ascoli tightness criterion therefore
gives tightness of the Rademacher restrictions to these rectangles.

Apply this argument on the exhaustion
$[j^{-1},j]$ and $[j^{-1},j]^2$, $j\ge2$.  The usual diagonal form of
the same criterion gives tightness in the product of the two
$C_{\mathrm{loc}}$ spaces.  Every subsequential limit has the Gaussian
finite-dimensional laws established above; evaluations on a countable
dense set determine a law on continuous functions.  The subsequential
limit is therefore $\Law_g(A_g,B_g)$, proving
Proposition~\ref{prop:universality}.

Every query in the comparison used the original times $sT,tT$ and the
original uniform initialization for $B_{n,J}$.  The replacement changes
only the coupling law.  Thus the coefficient $c_*$ and the entrance
$h_s^g$ obtained from the Gaussian construction pass unchanged to the
Rademacher limit.
